\scanpage{271}
\mspage{83}
% archive sheet number: 271
\unsure{redonnent} \inserted{(pour $\Phi\in\mathrm{Ob}\,\mathbb D(\unsure{K|k})$)} les flèches squelettiques

\[
\left\{\begin{aligned}
f^*(G)&\xrightarrow{u}F\\
G&\xrightarrow{v}f_*F
\end{aligned}\right.
\]

associées à $(F,G,u)$. De façon précise, $f^*(G)\longrightarrow F$ \deleted{s'identifie} se définit en termes de $\mathbb D$, et $\Phi\in\mathrm{Ob}\,\mathbb D(K)$, comme ci-dessus, et $v$ s'en déduit par adjonction.

L'idée naturelle ici, serait de ramener le pb de l'étoffement des flèches d'adjonction squelettiques, au cas des \unsure{imm.|immersions} ouvertes ou fermées. Notons que les sous-cat. ouvertes de $K$ correspondent justement aux foncteurs $K\longrightarrow\mathbb{\Delta}_{\unsure{9|1|q}}$, en prenant les images inverses de l'ouvert générique. Il y a donc un unique

\[
\unsure{\varphi|\phi}:K(f)\unsure{\text{---}|\longrightarrow}\mathbb{\Delta}_1
\]

avec $I=\varphi^{-1}(\{0\})$. Donc
